% !TEX root = handbook-physics.tex
%!TEX TS-program = pdflatex
%!TEX encoding = UTF-8 Unicode

\chapter{Ideal Gases}
\label{ch-08}
\zhtitle{理想气体}

\section{Subject Matter of Molecular Physics. Thermal Motion}
\label{sec-8.1}
\zhtitle{分子物理学的研究对象。热运动}
\subsection{}\label{8.1.1} The branch of physics called \redem{molecular physics} deals with the dependence of the structure and physical properties of bodies on the nature of the motion of and the interaction be­ tween the particles making up these bodies.
\zhpar{\redem{分子物理学}研究物体的结构与物理性质如何依赖于构成这些物体的粒子的运动性质及其相互作用。}

Molecular physics is based on the \redem{molecular-kinetic theory} of the structure of matter. According to this theory all bodies consist of submicroscopic particles -- atoms, molecules or ions -- in continuous chaotic motion, which is called \redem{thermal motion}. Experimental confirmations of the molecular-kinetic theory are: the Brownian movement \ssect{11.7.1}, the transport phenomenon in various aggregate states of matter and other phenomena. 
\zhpar{分子物理学以物质结构的\redem{分子动理论}为基础。根据该理论，一切物体都由亚微观粒子（原子、分子或离子）组成，这些粒子处于连续的无规则运动之中，这种运动称为\redem{热运动}。分子动理论的实验依据包括：布朗运动（见\ssect{11.7.1}）、物质各聚集态中的输运现象以及其他现象。}

\subsection{}\label{8.1.2} The molecular-kinetic theory of the structure of matter is used to advantage in various branches of physics. A great variety of physical phenomena that proceed due to the interac­tion and motion of particles of matter are dealt with in this theory from a unified point of view. This theory, for example, enables the mechanism of elasticity in solids \ssect{40.3.1} to be understood, reveals the cause of internal friction in gases \ssect{10.8.4} and liquids, explains the difference between real gases \ssect{12.1.2} and ideal gases \ssect{8.4.1}, and elucidates the mechanism of elec­trical conduction in conductors of electric current of various kinds, and the electrical and magnetic properties of matter.
\zhpar{物质结构的分子动理论在物理学的各个分支中得到了有效应用。由粒子的相互作用和运动所引起的各种物理现象，在该理论中被统一地加以研究。例如，该理论可以说明固体中弹性的机理（见\ssect{40.3.1}），揭示气体（见\ssect{10.8.4}）和液体内摩擦的成因，解释实际气体（见\ssect{12.1.2}）与理想气体（见\ssect{8.4.1}）之间的差别，并阐明各类导体中电传导的机理，以及物质的电学和磁学性质。}

\subsection{}\label{8.1.3} The thermal motion of particles of matter in various states of aggregation is not the same. It depends upon the forces of attraction or repulsion between the atoms, molecules and ions.
\zhpar{处于不同聚集态的物质粒子的热运动并不相同，它取决于原子、分子和离子之间的吸引力或排斥力。}

There are practically no attractive forces between the atoms and molecules in sufficiently dilute gases. This is due to the fact that the particles of such gases are located at distances from one another exceeding the range of molecular action \ssect{12.2.3}. The particles of such gases travel at uniform velocity and in a straight line until they collide with one another or with the walls of the container. These collisions are of a random nature. Each molecule in the volume of the gas may collide with any of the nearby particles, changing the direction of its motion in any arbitrary manner. The collision of a molecule (or atom) of the gas with a wall may occur at any angle $\theta$ where
$-\pi/2 \leqslant \theta \leqslant \pi/2$. As a result, the thermal motion of molecules of gas is a purely random one and, on the average, the same number of molecules are travelling at any instant of time in any arbitrary direction within the gas.
\zhpar{在足够稀薄的气体中，原子与分子之间几乎不存在吸引力。这是因为此类气体中粒子之间的距离超过了分子作用的范围（见\ssect{12.2.3}）。这类气体的粒子以匀速沿直线运动，直到它们彼此碰撞或与容器壁相碰为止。这些碰撞是随机的：气体中的每个分子都可能与附近的任一粒子发生碰撞，并以任意方式改变其运动方向。气体分子（或原子）与器壁碰撞时，碰撞角 $\theta$ 可取 $-\pi/2 \leqslant \theta \leqslant \pi/2$ 范围内的任意值。因此，气体分子的热运动是纯粹无规则的，平均而言，在任一时刻，气体内沿任意方向运动的分子数都相同。}

\subsection{}\label{8.1.4} Solid crystalline bodies are characterized by considerable forces of interaction between their particles (atoms, molecules and ions). The joint action of the forces of attraction and repul­sion between the particles \ssect{12.1.3} leads to the vibration of the particles of solids about their equilibrium positions, which are called the \redem{sites}, or \redem{points, of the crystal lattice} \ssect{40.1.1}. Intermolecular interaction and violations of periodicity in crystals (40.1.1) make these vibrations anharmonic \ssect{28.1.3}. 
\zhpar{固体晶体中粒子（原子、分子和离子）之间存在较强的相互作用力。粒子之间吸引力与排斥力的共同作用（见\ssect{12.1.3}）使固体粒子在其平衡位置附近振动，这些平衡位置称为\redem{格点}，或\redem{晶格结点}（见\ssect{40.1.1}）。分子间的相互作用以及晶体中周期性的破坏（40.1.1）使这些振动成为非简谐振动（见\ssect{28.1.3}）。}

\subsection{}\label{8.1.5} The thermal motion of molecules of a liquid is of a kind intermediate between the two preceding types of motion \ssect{8.1.3} and \ssect{8.1.4}. For a definite length of time the molecules of a liquid vibrate about a.certain equilibrium position and are in a tem­porarily settled state. At the end of this time the equilibrium position of the molecules is displaced and a new settled arrange­ ment is formed. Simultaneously, a slow displacement of the molecules and their vibrations occur within small volumes of the liquid (see also \ssect{13.2.3}).
\zhpar{液体中分子的热运动是介于上述两种运动形式之间的一种运动（见\ssect{8.1.3}和\ssect{8.1.4}）。在一定时间内，液体分子在某一平衡位置附近振动，处于暂时稳定的状态。经过这段时间后，分子的平衡位置发生移动，并形成新的稳定排列。与此同时，液体的小体积区域内还发生着分子的缓慢位移及其振动（亦见\ssect{13.2.3}）。}

\section[Statistical and Thermodynamic Methods]{Statistical and Thermodynamic\\ Methods of Investigation}
\label{sec-8.2} 
\zhtitle{统计方法与热力学研究方法}

\subsection{}\label{8.2.1} There is a huge number of atoms (or molecules) in any body. For instance, in \SI{1}{\centi\meter\cubed} of a gas whose properties are close to those of an ideal gas there are \num{2.7e19} molecules at stand­ard conditions. In the condensed states, liquid and solid, the amount is of the order of \num{d22} particles per \si{\centi\meter\cubed}. If we assume that each atom (or molecule) obeys Newton's second law \ssect{2.4.1} it is obviously quite impossible, not only to solve the dif­ferential equations of motion of the separate particles, but even to write these equations. Therefore, the behaviour of a separate molecule (or atom) of a body, for example, its path and the sequence of changes in its state, cannot be investigated by the methods used in classical mechanics.
\zhpar{任何物体中都含有数量极其庞大的原子（或分子）。例如，在标准状态下，性质接近理想气体的气体，每 \SI{1}{\centi\meter\cubed} 中含有 \num{2.7e19} 个分子。在凝聚态（液态和固态）中，每 \si{\centi\meter\cubed} 的粒子数量级为 \num{d22} 个。若假定每个原子（或分子）都服从牛顿第二定律（见\ssect{2.4.1}），则显然不仅无法求解每个粒子运动的微分方程，甚至连写出这些方程都不可能。因此，物体中单个分子（或原子）的行为，例如其轨迹以及状态变化的顺序，不能用经典力学的方法加以研究。}

\subsection{}\label{8.2.2} The macroscopic properties of systems consisting of a very great number of particles are studied by the \redem{statistical method}. It is based on the application of probability theory to definite models of the structure of the systems being investigated. The branch of theoretical physics in which the physical prop­erties of systems are studied by the statistical method is called \redem{statistical physics (statistical mechanics)}. Special laws, known as statistical laws are manifested in the joint behaviour of large numbers of particles. Certain \redem{average values of the physical quantities} can be found for a system consisting of a great many particles. These values characterize the set of particles as a whole. In gases, for instance, there are average values of the velocity of thermal motion of the molecules \ssect{10.3.6} and of their energy \ssect{10.2.4}. In a solid there is an average energy per degree of freedom of the vibrational motion of its particles \ssect{41.7.2}, etc. All properties of a system of particles are due, not only to the individual properties of the particles themselves, but to the features of their joint motion and average values of their dynamic characteristics (average velocities, average energies, etc.).
\zhpar{由大量粒子组成的系统的宏观性质，用\redem{统计方法}进行研究。该方法以概率论为基础，并应用于所研究系统结构的确定模型。用统计方法研究系统物理性质的理论物理分支，称为\redem{统计物理学（统计力学）}。大量粒子的联合行为中会表现出一类特殊的规律，称为统计规律。对于由大量粒子组成的系统，可以求得某些\redem{物理量的平均值}，这些平均值表征整个粒子集合的特性。例如，气体中存在分子热运动速度的平均值（见\ssect{10.3.6}）以及分子能量的平均值（见\ssect{10.2.4}）。固体中粒子振动运动每个自由度的平均能量（见\ssect{41.7.2}）等也属此类。粒子系统的所有性质，不仅取决于粒子本身的个体性质，还取决于它们的联合运动特征以及其动力学特征的平均值（平均速度、平均能量等）。}

In addition to statistical laws, there are also dynamical laws, describing the motion of separate particles. The relationship between dynamical and statistical laws is manifested in the fact that the laws of motion of individual particles have an effect of the description of a system of particles, which is being investigated by the statistical method.
\zhpar{除统计规律外，还存在描述单个粒子运动的动力学规律。动力学规律与统计规律之间的关系体现在：单个粒子的运动规律会影响用统计方法所研究的粒子系统的描述。}

\subsection{}\label{8.2.3} Another procedure for studying physical phenomena is the \redem{thermodynamic method}, which is not concerned with either the internal structure of the substances making up the bodies (or systems) being investigated or with the nature of the motion of the separate particles. The thermodynamic method is based on the study of the various conversions of energy that occur in a system. The conditions of these conver­sions and the relations between the various kinds of energy enable the physical properties of systems to be studied in a great variety of processes in which the systems participate. The branch of physics in which research in the physical properties of systems is conducted by means of the thermodynamic method is called \redem{thermodynamics (phenomenological thermody­namics)}. Thermodynamics is based on two experimentally established principles, or laws, of thermodynamics \ssect{9.3.1} and \ssect{11.3.2}, as well as the Nernst. heat theorem, or third law of thermodynamics \ssect{11.8.2}.
\zhpar{研究物理现象的另一种方法是\redem{热力学方法}。它既不涉及所研究物体（或系统）的物质内部结构，也不涉及单个粒子的运动性质。热力学方法基于对系统中各种能量转化过程的研究。这些转化的条件以及各种能量之间的关系，使人们能够在系统参与的各种过程中研究系统的物理性质。以热力学方法研究系统物理性质的物理学分支，称为\redem{热力学（唯象热力学）}。热力学以两条由实验确立的原理（即定律）为基础，即热力学第一定律（见\ssect{9.3.1}）和热力学第二定律（见\ssect{11.3.2}），此外还有能斯特热定理，即热力学第三定律（见\ssect{11.8.2}）。}

\clearpage

\section[Thermodynamic Variables]{Thermodynamic Variables. Equations of State. Thermodynamic Processes}
\label{sec-8.3}
\zhtitle{热力学变量。物态方程。热力学过程}

\subsection{}\label{8.3.1} Thermodynamics deals with \redem{thermodynamic systems}. There are macroscopic items (bodies and fields) that can interchange energy with one another as well as with the external medium, i.e. bodies and fields that are external with respect to the given system.
\zhpar{热力学研究\redem{热力学系统}。系统之外存在一些宏观的物体和场，它们可以与该系统以及外部介质交换能量，即相对于所给定的系统而言处于外部的物体和场。}

\subsection{}\label{8.3.2} The physical quantities introduced to describe the state of a thermodynamic system are called \redem{thermodynamic variables}, or \redem{parameters (state variables)}, of the system. Usually the pres­sure, specific volume and temperature are taken as the thermo­ dynamic variables.
\zhpar{用来描述热力学系统状态的物理量，称为该系统的\redem{热力学变量}，或\redem{参量（状态参量）}。通常取压强、比体积和温度作为热力学变量。}

The \redem{pressure} $p$ is a physical quantity equal to the force acting on unit area of a body's surface in the direction normal to the surface:
\zhpar{\redem{压强} $p$ 是一个物理量，它等于作用在物体表面单位面积上、沿表面法线方向的力，其定义式为：}
\begin{equation*}%
p = \dv{F_{n}}{A},
\end{equation*}
where $\dd F_{n}$ is the numerical value of the normal force acting on an infinitely small part of the surface of the body, the area of this part being $\dd A$.
\zhpar{式中 $\dd F_{n}$ 为作用在物体表面某一无穷小部分上的法向力的数值，该部分的面积为 $\dd A$。}

The \redem{specific volume} $v$ is the reciprocal of the density $\rho$ of a body: $v = 1/\rho$. The specific volume of a homogeneous body is equal to the volume of unit mass of the body.
\zhpar{\redem{比体积} $v$ 是物体密度 $\rho$ 的倒数：$v = 1/\rho$。均匀物体的比体积等于该物体单位质量的体积。}

\subsection{}\label{8.3.3} The concept of temperature is of significance for equilibrium states of a thermodynamic system \ssect{8.3.4}. An \redem{equilib­rium state (state of thermodynamic equilibrium)} is a state of a system which does not change with time (\redem{steady state}), the steady nature of the state not being due to processes occurring in the external medium. An equilibrium state is set up in a system subject to constant external conditions and is main­tained in the system for an arbitrary length of time. The tem­perature is the same in all parts of a thermodynamic system that is in a state of thermodynamic equilibrium. If two bodies with different temperatures come into contact, internal ener­gy \ssect{9.1.2} is transferred from the body at the higher temperature to that at the lower temperature. The process ceases when the temperatures of the two bodies become equal.
\zhpar{温度的概念对热力学系统的平衡态具有重要意义（见\ssect{8.3.4}）。\redem{平衡态（热力学平衡态）}是指系统的状态不随时间变化（\redem{稳定态}）的状态，且该状态的稳定性并非由外部介质中的过程引起。在恒定的外部条件下，系统中会建立平衡态，并能在任意长的时间内保持。处于热力学平衡态的系统，其各部分温度都相同。若两个温度不同的物体相互接触，内能（见\ssect{9.1.2}）将从高温物体传递到低温物体。当两物体的温度相等时，该过程停止。}

\subsection{}\label{8.3.4} The temperature of an equilibrium system is a measure of the intensity of thermal motion of its molecules (or atoms or ions). In an equilibrium system of particles that obeys the laws of classical statistical mechanics \ssect{8.2.2}, the average kinetic energy of the thermal motion of the particles is pro­portional to the absolute temperature \ssect{10.2.5}.
\zhpar{平衡系统的温度是其分子（或原子、离子）热运动剧烈程度的量度。对于服从经典统计力学规律的平衡粒子系统（见\ssect{8.2.2}），粒子热运动的平均动能与绝对温度成正比（见\ssect{10.2.5}）。}

Temperature can only be measured indirectly, making use of the fact that many physical properties of bodies, lending themselves to direct measurement, depend on the temperature. A change in the temperature of a body leads to changes in its length and volume, density, electric resistance, elastic prop­erties, etc. A change in any of these properties can be used to measure temperature. For this purpose, it is necessary for a certain body, said to be \redem{temperature-indicating}, that the functional dependence of the given property on the temperature be known. Temperature scales based on temperature-indicating bodies are said to be \redem{empirical}.
\zhpar{温度只能间接测量，这是利用了物体的许多可直接测量的物理性质依赖于温度这一事实。物体温度的变化会引起其长度、体积、密度、电阻、弹性性质等的变化。利用这些性质中的任何一种变化都可以测量温度。为此，必须已知某一称为\redem{测温物质}的物体的上述性质与温度之间的函数关系。以测温物质为基础建立的温标称为\redem{经验温标}。}

In the \redem{international centigrade (Celsius)} scale, temperature is measured in  \si{\degreeCelsius} and is denoted by $t$ (degree centigrade or degree Celsius). It is assumed that at a standard pressure of \SI{1.01325d5}{\pascal} (Appendix~I) the melting point of ice and the boiling point of water are equal to \SI{0}{\degreeCelsius} and \SI{100}{\degreeCelsius}.
\zhpar{在\redem{国际摄氏温标}中，温度以 \si{\degreeCelsius} 为单位，记为 $t$（摄氏度）。假定在标准大气压 \SI{1.01325d5}{\pascal}（见附录I）下，冰的熔点和水的沸点分别为 \SI{0}{\degreeCelsius} 和 \SI{100}{\degreeCelsius}。}

In the thermodynamic (absolute) temperature scale, the tem­perature is measured in kelvins (\si{\kelvin}) and is denoted by $T$.
\zhpar{在热力学（绝对）温标中，温度以开尔文（\si{\kelvin}）为单位，记为 $T$。}

The relation between the absolute temperature and the tem­perature measured by the Celsius scale is: 
\zhpar{绝对温度与摄氏温标所测温度之间的关系如下式所示：}
\begin{equation*}%
T = 273.15 + t.
\end{equation*}
The temperature $T = \SI{0}{\kelvin}$ (or $t = \SI{-273.15}{\degreeCelsius}$) is called absolute zero temperature (the impossibility of reaching absolute zero is discussed in \ssect{11.8.4}).
\zhpar{温度 $T = \SI{0}{\kelvin}$（或 $t = \SI{-273.15}{\degreeCelsius}$）称为绝对零度（关于无法达到绝对零度的讨论见\ssect{11.8.4}）。}

\subsection{}\label{8.3.5} The state variables of a system are classified as being either external or internal. \redem{External state variables} are physical quantities depending upon the position in space and various properties of bodies that are external with respect to the given system. For a gas, for instance, the volume $V$ of the vessel in which the gas is contained is an external variable because the volume depends upon the position of external bodies: the walls of the vessel. For a dielectric in an electric field, the intensity of the field, set up by certain external sources, is an external variable. If a liquid is held in an open vessel, the atmospheric pressure, for example, is an external variable. 
\zhpar{系统的状态参量可分为外部参量和内部参量两类。\redem{外部状态参量}是取决于空间位置以及相对于该系统处于外部的物体性质的物理量。例如，对于气体，盛放气体的容器体积 $V$ 是外部参量，因为其体积取决于外部物体（即器壁）的位置。对于处于电场中的电介质，由某些外部源建立的电场强度是外部参量。若液体盛放在敞口容器中，例如大气压强就是外部参量。}

The \redem{internal variables of a system} are physical quantities that depend both on the position of bodies external with respect to the system, and on the coordinates and velocities of the particles that make up the system. Internal variables of a gas, for instance, are its pressure and energy because they depend upon the coordinates and velocities of the moving molecules and on the density of the gas.
\zhpar{\redem{系统的内部参量}是既取决于相对于该系统处于外部的物体的位置，又取决于构成该系统的粒子的坐标和速度的物理量。例如，气体的内部参量是其压强和能量，因为它们取决于运动分子的坐标和速度以及气体的密度。}

\subsection{}\label{8.3.6} The state variables of a system that is in equilibrium \ssect{8.3.3} are not independent. The internal variables of such a system depend only on its external variables and on the temperature. The equilibrium state of a simple system\footnote{An example of a simple system is a gas in the absence of external fields. Mixtures of chemically homogeneous gases of constant composition, chemically pure liquids, etc. arc also simple systems.} of given chemical composition and mass $M$ is determined by specifying two variables: the volume $V$ and the temperature $T$. The equation of state (thermal equation of state) of a simple system is the functional dependence of the equilibrium pressure $p$ in the system on the volume and temperature
\zhpar{处于平衡态的系统（见\ssect{8.3.3}）的状态参量并不独立。此类系统的内部参量仅取决于其外部参量和温度。化学组成给定、质量为 $M$ 的简单系统的平衡态由两个参量确定：体积 $V$ 和温度 $T$。简单系统的物态方程（热物态方程）是其平衡压强 $p$ 与体积和温度的函数关系，即：}
\zhpar{注：简单系统的一个例子是不存在外场的气体。化学成分均匀且组成恒定的混合气体、化学纯液体等也都是简单系统。}
\begin{equation*}%
p= f(V, T).
\end{equation*}
The equation of state is obtained experimentally in thermo­ dynamics. In statistical physics \ssect{8.2.2}, it is derived theoret­ically. This constitutes the interrelation between the statis­tical \ssect{8.2.2} and thermodynamical \ssect{8.2.3} methods of investi­gation.
\zhpar{物态方程在热力学中通过实验获得，而在统计物理学中则由理论推导得出。这构成了统计方法（见\ssect{8.2.2}）与热力学方法（见\ssect{8.2.3}）之间的内在联系。}

\subsection{}\label{8.3.7} If any of the external variables of the system is changed, the state of a thermodynamic system is changed. This is called a \redem{thermodynamic process}. A thermodynamic process is said lo be equilibrium (an \redem{equilibrium, quasi-static process}) if the system passes at an infinitely slow rate through a continuous succession of infinitely close thermodynamic equilibrium states \ssect{8.3.3}. All processes that do not comply with these conditions are said to be \redem{non-equilibrium} ones. Real processes are non-equilibrium ones because they proceed at a finite rate. But the slower they proceed, the closer they are to equilibrium processes
\zhpar{若系统的任一外部参量发生变化，热力学系统的状态也随之改变，这称为\redem{热力学过程}。若系统以无限缓慢的速率经历一系列无限接近的热力学平衡态，则该热力学过程称为平衡过程（\redem{平衡过程}，或\redem{准静态过程}）。不满足上述条件的过程称为\redem{非平衡}过程。实际过程都是非平衡过程，因为它们以有限速率进行，但进行得越慢，就越接近平衡过程。}

\redem{Isoprocesses} are thermodynamic processes that proceed in a system of constant mass with a constant value of one of the state variables.
\zhpar{\redem{等值过程}是指在质量不变的系统中，某一状态参量保持恒定时所进行的热力学过程。}

An \redem{isothermal process} proceeds at constant temperature $(T = \textrm{const})$.
\zhpar{\redem{等温过程}是在恒定温度 $(T = \textrm{const})$ 下进行的过程。}

An \redem{isochoric (isovolumic) process} proceeds at constant volume $(V = \textrm{const})$.
\zhpar{\redem{等容（等体）过程}是在恒定体积 $(V = \textrm{const})$ 下进行的过程。}

An \redem{isobaric process} proceeds at constant pressure $(p = \textrm{const})$. 
\zhpar{\redem{等压过程}是在恒定压强 $(p = \textrm{const})$ 下进行的过程。}

An \redem{adiabatic process} is a thermodynamic process that occurs in a system without heat exchange \ssect{9.2.4} with external bodies (see also \ssect{9.5.10}).
\zhpar{\redem{绝热过程}是指系统与外部物体之间没有热量交换（见\ssect{9.2.4}）的热力学过程（亦见\ssect{9.5.10}）。}

\subsection{}\label{8.3.8} \redem{State functions} are physical quantities characterizing the state of a system. Variation in the state functions during ther­modynamic processes is independent of the kind of process that occurs. State functions are uniquely determined by the values of the variables of the initial and final states of the system. The simplest state functions of a system are its inter­nal energy $U$ \ssect{9.1.2} and its entropy $S$ \ssect{11.4.2}.
\zhpar{\redem{状态函数}是表征系统状态的物理量。热力学过程中状态函数的变化与过程的类型无关。状态函数完全由系统初态和终态的参量值唯一确定。系统最简单的状态函数是其内能 $U$（见\ssect{9.1.2}）和熵 $S$（见\ssect{11.4.2}）。}

\section{Equation of State of an Ideal Gas}
\label{sec-8.4}

\subsection{}\label{8.4.1} An \redem{ideal}, or \redem{perfect, gas} is one whose molecules do not interact with one another at a distance and are of vanishingly small size. In real gases \ssect{12.1.2} the molecules are subject to the forces of intermolecular interaction \ssect{12.1.3}.
\zhpar{\redem{理想气体}，或称\redem{完全气体}，是指其分子之间在一定距离上不存在相互作用、且线度可以忽略为零的气体。在实际气体中（见\ssect{12.1.2}），分子受到分子间相互作用力（见\ssect{12.1.3}）的作用。}

In collisions with one another and in striking the walls of the vessel the molecules of an ideal gas behave as perfectly elastic spheres \ssect{1.1.4} of diameter $d$ (\redem{elective diameter of the molecule}), which depends upon the chemical nature of the gas. The existence of an effective diameter $d$ (where $d \approx \SI{d11}{\meter}$) signifies that forces of mutual repulsion \ssect{12.1.3} act between the molecules. Intermolecular forces of attraction \ssect{12.1.3} decrease drastically with an increase in the distance $r$ between molecules and are practically negligible at $ R_{m} > \SI{d9}{\meter}$. In real gases at low densities the average distance $\expval{r}$) between molecules exceeds $R_{m}$ and, to a good approximation, they can be regarded as ideal. Hydrogen, helium, oxygen and nitrogen are assumed to be ideal gases at densities corresponding to standard conditions.
\zhpar{在相互碰撞以及与容器壁相撞时，理想气体分子表现为直径为 $d$ 的完全弹性球（见\ssect{1.1.4}）。$d$ 称为\redem{分子的有效直径}，其数值取决于气体的化学性质。有效直径 $d$（其中 $d \approx \SI{d11}{\meter}$）的存在表明分子之间存在相互排斥力（见\ssect{12.1.3}）。分子间的吸引力（见\ssect{12.1.3}）随分子间距离 $r$ 的增大而迅速减小，当 $ R_{m} > \SI{d9}{\meter}$ 时实际上可以忽略。在低密度的实际气体中，分子间的平均距离 $\expval{r}$ 超过 $R_{m}$，在较好的近似下可以将它们视为理想气体。在标准状态下，氢、氦、氧和氮的密度都可以视为理想气体的密度。}

\subsection{}\label{8.4.2} For a given mass of an ideal gas, the ratio of the product of the numerical values of pressure and volume to the absolute temperature is a constant value (\redem{Clapeyron equation}):
\zhpar{对于给定质量的理想气体，压强与体积的数值乘积和绝对温度之比为常量（\redem{克拉珀龙方程}）：}
\begin{equation*}%
\frac{pV}{T} = C = \textrm{const.}
\end{equation*}
The numerical value of the \redem{gas constant} $C$ depends upon the choice of the units of measurement of $p$, $V$ and $T$.
\zhpar{\redem{气体常量} $C$ 的数值取决于 $p$、$V$ 和 $T$ 所采用的测量单位的选择。}

If $v$ is the specific volume of the gas \ssect{8.3.2} and $M$ is its mass, then $V = Mv$ and the Clapeyron equation assumes the form
\zhpar{若 $v$ 为气体的比体积（见\ssect{8.3.2}），$M$ 为其质量，则 $V = Mv$，克拉珀龙方程取如下形式：}
\begin{equation*}%
pv = \frac{C}{M} T = BT.
\end{equation*}
where $B = C/M$ is the \redem{specific gas constant} referred to unit mass.
\zhpar{式中 $B = C/M$ 为单位质量气体的\redem{比气体常量}。}

\subsection{}\label{8.4.3} From the definition of the mole (Appendix~I), it is evident that moles of different gases contain the same number of mole­cules. This is called \redem{Avogadro's number} $N_{A}$ (Appendix~II).
\zhpar{由摩尔的定义（见附录I）可知，不同气体的摩尔所含分子数相同。该数称为\redem{阿伏伽德罗常量} $N_{A}$（见附录II）。}

The \redem{molar mass} $\mu$ of a gas or any body is the physical quantity equal to the ratio of the mass of the gas (or any body) to the number of moles $N$ it contains: $\mu = M/N$.
\zhpar{气体或任何物体的\redem{摩尔质量} $\mu$ 是其质量与所含摩尔数 $N$ 之比：$\mu = M/N$。}

The molar mass is proportional to the relative mass of the molecules of a gas: $\mu = \num{d-3} m/m_{0}$, where $m$ is the mass of the molecules of the given gas and $m_{0}$ is the atomic mass unit based on the carbon scale (Appendix~II). The coefficient \num{d-3} is required because molar mass is measured in \si{\kilo\gram\per\mole} in the SI (International System of Units).
\zhpar{摩尔质量与气体分子的相对质量成正比：$\mu = \num{d-3} m/m_{0}$，其中 $m$ 为给定气体分子的质量，$m_{0}$ 为以碳标度为基准的原子质量单位（见附录II）。系数 \num{d-3} 是必需的，因为在国际单位制（SI）中，摩尔质量以 \si{\kilo\gram\per\mole} 为单位。}

The \redem{molar volume} $V_{\mu}$ is the physical quantity equal to the ratio of the volume $V$ of a gas to the number of moles $N$ it contains: $V_{\mu} = V/N$. The mass of a mole is numerically equal to $\mu$ and therefore $V_{\mu} = \mu v$, where $v$ is the specific volume \ssect{8.3.2}.
\zhpar{\redem{摩尔体积} $V_{\mu}$ 是气体的体积 $V$ 与其所含摩尔数 $N$ 之比：$V_{\mu} = V/N$。一摩尔气体的质量在数值上等于 $\mu$，因此 $V_{\mu} = \mu v$，其中 $v$ 为比体积（见\ssect{8.3.2}）。}

\subsection{}\label{8.4.4} The equation of state for a mole of ideal gas is
\zhpar{一摩尔理想气体的物态方程为}
\begin{equation*}%
\frac{pV_{\mu}}{T} = \mu B = R \qor pV_{\mu} = RT.
\end{equation*}
Here $R$ is the \redem{universal gas constant}, being the gas constant per mole of gas. The universal nature of $R$ follows from \redem{Avagadro's law}, which states that moles of all ideal gases occupy the same volume at the same pressure and the same tempera­ture. Under standard conditions ($T = \SI{273.15}{\kelvin} and $p =\SI{1.0132d5}{\pascal} = \SI{1}{\atmos} = \SI{760}{\milli\meter\mercury}) a mole of any gas has a volume of $V_{\mu} = \SI{22.415d-3}{\meter\cubed}$. From this we can calculate the numerical values of $R$ for various systems of units (Appendix~II).
\zhpar{其中 $R$ 为\redem{摩尔气体常量}，即每摩尔气体的气体常量。$R$ 的普适性源于\redem{阿伏伽德罗定律}：在相同的压强和温度下，任何理想气体的一摩尔所占体积都相同。在标准状态（$T = \SI{273.15}{\kelvin}$，$p = \SI{1.0132d5}{\pascal} = \SI{1}{\atmos} = \SI{760}{\milli\meter\mercury}$）下，任何气体一摩尔的体积为 $V_{\mu} = \SI{22.415d-3}{\meter\cubed}$。由此可以算出 $R$ 在各种单位制中的数值（见附录II）。}

If a volume $V$ of gas has the mass $M$ \si{\kilo\gram}, i.e. if it contains $M/\mu$ moles, then $V = M /\mu \times V_{\mu}$ and the equation of state of the gas assumes the form called the \redem{Mendeleev-Clapeyron equation}:
\zhpar{若气体体积 $V$ 的质量为 $M$ \si{\kilo\gram}，即其含有 $M/\mu$ 摩尔，则 $V = M /\mu \times V_{\mu}$，气体的物态方程取如下形式，称为\redem{门捷列夫-克拉珀龙方程}：}
\begin{equation*}%
pV = \frac{M}{\mu} RT.
\end{equation*}

\subsection{}\label{8.4.5} \redem{Boltzmann's constant} k is a physical quantity equal to the ratio of the universal gas constant $R$ to Avogadro's number $N_{A}$, i.e. $k = R/N_{a}$. The values of $k$ in various systems of units are listed in Appendix~II.
\zhpar{\redem{玻尔兹曼常量} $k$ 是一个物理量，它等于摩尔气体常量 $R$ 与阿伏伽德罗常量 $N_{A}$ 之比，即 $k = R/N_{A}$。$k$ 在各种单位制中的数值列于附录II中。}

The equation of state of an ideal gas, expressed by means of Boltzmann's constant, is of the form
\zhpar{用玻尔兹曼常量表示的理想气体物态方程形式为}
\begin{equation*}%
p = \frac{kN_{A}}{V_{\mu}} = kn_{0}T,
\end{equation*}
where $n_{0}= N_{A}/V_{\mu}$ is the number of gas molecules in unit volume (molecule concentration).
\zhpar{式中 $n_{0}= N_{A}/V_{\mu}$ 为单位体积内的气体分子数（分子浓度）。}

At constant temperature, the pressure of a gas is proportional to the concentration of its molecules (or to the density of the gas).
\zhpar{在温度恒定时，气体的压强与其分子浓度（或气体密度）成正比。}

